\section{Background and Motivation}
\label{sec:background}

\subsection{Field-Level Conformance}

Protocol conformance asks whether an implementation satisfies the behavioral
requirements stated by a standard. For natural-language standards, this is a
specification-to-code alignment problem: the checker must connect a normative
requirement to the protocol object it constrains and to the implementation
behavior that determines whether the requirement is satisfied. Incorrect
alignment can produce a wrong conclusion, for example when a length rule is
matched to an unrelated buffer or when a validation path is missed because it is
enforced outside the parser.

This work focuses on \emph{field-level semantic conformance}. Fields are a
useful boundary because standards define them through message formats,
parameter tables, and normative prose, while implementations eventually
represent them through parsed values, state variables, or validation logic. The
names may differ, but the field gives the checker a stable object to track.
This scope is narrower than full protocol verification, but it covers many
practical defects because parsing, validation, and error handling are usually
organized around fields.

\subsection{Normative Modalities}

The strength of a conformance claim depends on normative modality. Requirements
phrased with \emph{MUST} or \emph{MUST NOT} usually support direct violation
reports. Requirements phrased with \emph{SHOULD} are tracked but reported as
lower-severity deviations, while \emph{MAY} usually indicates optional behavior
and is ignored as a primary check target. \tool{} records modality as part of
each rule and uses it to rank report severity. A \emph{MUST NOT accept} rule
with a reproducible accepting input is reported as a strong inconsistency. A
\emph{SHOULD reject} rule with similar evidence is reported as a SHOULD-level
deviation rather than a low-confidence result. This distinction is important
because protocol standards often encode both hard interoperability constraints
and softer deployment guidance.

\subsection{A Running Example}

Consider a simplified QUIC transport rule: if an endpoint receives a transport
parameter that violates a specified bound, it must terminate the connection with
a protocol error. The standard may express this requirement across a field
definition, a paragraph about transport parameters, and an error-handling
section. The implementation may parse the parameter in one file, store it in a
connection configuration structure, validate it in a helper, and trigger the
error through a callback. A direct text-to-code query such as ``does the code
enforce the transport parameter bound?'' is under-specified: the relevant
variable names, call chain, and build configuration are all missing.

\tool{} handles the example by first extracting a rule such as
\[
\begin{aligned}
C &= \textit{parameter received},\\
f &= \textit{max\_udp\_payload\_size},\\
A &= \textit{reject values below the lower bound}.
\end{aligned}
\]
It then searches for identifiers and code regions related to the field,
including aliases such as \texttt{max\_udp\_payload},
\texttt{udp\_payload\_size}, or structure members embedded in parser state.
If the first retrieved context only shows assignment but not validation, the
reasoning agent emits \emph{insufficient evidence} and requests callers or
validation helpers. Only after the evidence covers the relevant path does the
system report a consistency judgment. For a suspected deviation, the test agent
attempts to construct a packet or configuration that exercises the suspicious
path and observes whether the implementation accepts or rejects it.

% \subsection{Threats to Naive LLM Checking}

% LLMs are useful because they can relate prose to code-level intent, but
% conformance checking punishes ungrounded flexibility. First, field names are
% ambiguous: a standard may use a short name in one paragraph and a fully
% qualified message path in another. Second, code identifiers are indirect:
% \texttt{len}, \texttt{payload\_len}, and \texttt{param->value} may refer to
% different protocol objects depending on scope. Third, evidence is often
% negative: a missing check is meaningful only after the checker has searched
% the right parser, validation helper, and caller constraints. Finally, execution
% semantics may depend on macros, compile-time options, and runtime configuration.
% \tool{} is designed around these threats. It constrains field references,
% retrieves code structurally, records uncertainty, and validates high-risk
% claims dynamically when possible.
